\chapter{Maneuver-Level Signal Understanding: An In-Depth Analysis}
\label{ch:problem}

This chapter examines the task itself before any model is trained. It gives a formal definition
(\cref{sec:prob-def}), explains how the abstract notion of ``the signal that governs a maneuver'' is
made concrete in labels and why this depends on traffic rules (\cref{sec:prob-semantics}), discusses
the characteristics of the data that make the task hard (\cref{sec:prob-challenges}), analyses the
evaluation metrics and what they do not show (\cref{sec:prob-metrics}), and contrasts the two model
designs compared in this thesis (\cref{sec:prob-paradigms}).

\section{Task definition}
\label{sec:prob-def}

Let $I$ be a front-camera image and $\dirs=\{\text{straight},\text{left},\text{right}\}$
the set of maneuvers. The task is to predict, for every $d\in\dirs$,
\begin{align}
  c_d &\in \mathcal{C} = \{\text{none},\text{red},\text{yellow},\text{green}\}, \\
  p_d &\in \mathcal{P} = \{\text{stop},\text{go}\}.
\end{align}
The color $c_d$ is the state of the signal that governs maneuver $d$ for the ego vehicle at the next
intersection, or \emph{none} when no signal governs it. The passability is defined by the color through
the rule of Zou et al.\ (legend of their Fig.~2):
\begin{equation}
  p_d = \begin{cases} \text{stop} & \text{if } c_d = \text{red},\\ \text{go} & \text{otherwise.}\end{cases}
  \label{eq:pass-rule}
\end{equation}
Yellow therefore counts as \emph{go}, and so does \emph{none}: a maneuver without a governing signal is
not prohibited by a signal, even if other rules (yield, give way) apply. The model still predicts $p_d$
with a separate head, which the original paper motivates with the possibility of learning exceptions;
\cref{sec:res-perdir} shows that in practice the two outputs are almost perfectly coupled.

A frame-level label is the tuple $y = \big((c_d, p_d)\big)_{d\in\dirs}$. Since $p_d$ is a function of
$c_d$, the label space has $4^3=64$ possible states, of which far fewer occur in practice because
maneuvers are strongly correlated: straight and left share a signal whenever there is no dedicated left
arrow.

\paragraph{Why U-turn is left out.} Zou et al.\ define a fourth maneuver, the U-turn. This thesis drops
it. Dedicated U-turn signals are almost absent from the data, so a U-turn label takes the left signal in
99.9\% of the OpenLane-V2 frames, and DTLD and LISA have no U-turn signals at all. As a separate output,
U-turn therefore adds no information; it only counts the left turn a second time in every averaged
metric. All results in this thesis are computed over the three maneuvers, unless stated otherwise. The
paper's published numbers include U-turn, so comparisons with them are close but not exact. Removing
a near copy of left reweights the averaged metrics towards straight and right; since right is almost
always the trivial \emph{none} (\cref{sec:prob-metrics}), three-direction scores come out slightly
higher than four-direction scores of the same predictions.

\section{Label semantics and traffic rules}
\label{sec:prob-semantics}

The definition above hides the hardest part of the problem: deciding \emph{which} light governs a
maneuver. In OpenLane-V2 this has to be derived from the lane graph, and several choices have to be
made. The rules adopted in this thesis, detailed in \cref{sec:method-labels} and
\cref{app:labels}, follow a hierarchy of specificity:
\begin{enumerate}
  \item \textbf{Straight} takes the color of the lights that control the straight lane or that are
        generic (circular) lights of the approach.
  \item \textbf{Left} takes a left-specific light if one exists, and otherwise inherits the straight
        signal.
  \item \textbf{Right} takes only a \emph{dedicated} right-turn light; without one it is \emph{none}.
\end{enumerate}
The third rule encodes the US convention that a right turn on red is permitted
(\cref{sec:bg-signals}). It has two consequences that recur throughout the thesis. First, right is
\emph{none} in about 97\% of the OpenLane-V2 frames, so its color accuracy is dominated by a trivial
majority class. Second, the rare dedicated red right arrow, which is the one case where turning right
is actually forbidden, is extremely under-represented.

The convention also does not transfer. Under German rules (DTLD) the right turn obeys the circle when
there is no right arrow; applying the US rule there would declare most right turns unconstrained and
would be wrong. The DTLD labels in this thesis therefore use a second convention (\emph{signal}: right
takes a right arrow if present, otherwise the circle) and record it per frame, so that both training
and evaluation apply the correct rule. Right-turn labels thus mean different things across datasets,
which matters when models are evaluated out of distribution (\cref{sec:res-ood}).

A further subtlety concerns the left-hand traffic scenes of nuScenes (Singapore). A strict mirroring of
the rules (left becomes the near turn, right the far turn) would be the natural choice, but, as
shown in \cref{sec:res-labels}, it does not reproduce the label distribution of the original paper,
whereas applying US rules everywhere does. The replication therefore follows the US convention for
comparability, at the price of a known inconsistency: in Singapore a green right arrow yields
right~$=$~none.

\section{Characteristics and challenges of the data}
\label{sec:prob-challenges}

\paragraph{Small objects.} Traffic lights at the stop line of the next intersection are often 10--30
pixels tall in a 2048-pixel image, and they shrink by a factor of two to three when the image is
resized to a size a network can process. \Cref{sec:res-atlas} quantifies how strongly detection
recall depends on this height.

\paragraph{Relevance.} A frame can contain lights of the ego lane, of neighboring lanes with different
maneuvers, of the opposite approach (seen from the side), and of the next intersection. Lights of a far
intersection can have a different color from the near ones, and the label builder has to decide which
intersection is ``next''. When the most specific relevant lights disagree in color, the situation is
ambiguous and such frames are excluded from the benchmark.

\paragraph{Class imbalance.} Beyond the right-turn imbalance, yellow is rare (about 3\% of straight
labels) and short-lived, and \emph{none} for straight (T-junctions, turn-only lanes) is rare as well.
Rare classes carry little weight in pooled accuracy but include some of the most consequential
decisions.

\paragraph{Correlated frames.} The datasets consist of driving sequences sampled at 2\,Hz or more.
Consecutive frames usually show the same lights in the same state, so the effective number of
independent samples is closer to the number of intersection approaches than to the number of frames.
This has to be taken into account when comparing models statistically (\cref{sec:method-eval}).

\paragraph{Domain shift.} Lights differ in housing, size, mounting (overhead versus side poles) and
color rendition between countries and cameras; night images change their appearance completely; and
the camera geometry (portrait Argoverse~2 images, wide DTLD images, 4:3 LISA images) changes the scale
and aspect ratio at which lights appear after resizing.

\section{Evaluation metrics and their blind spots}
\label{sec:prob-metrics}

The metrics follow Eqs.~4--7 of Zou et al. For $N$ frames with predictions $\hat c_{i,d}$, $\hat
p_{i,d}$:
\begin{align}
  \acccolor &= \frac{1}{|\dirs|N}\sum_{i=1}^{N}\sum_{d\in\dirs} \mathbb{1}[\hat c_{i,d} = c_{i,d}], \\
  \accpass &= \frac{1}{|\dirs|N}\sum_{i=1}^{N}\sum_{d\in\dirs} \mathbb{1}[\hat p_{i,d} = p_{i,d}], \\
  \EM &= \frac{1}{N}\sum_{i=1}^{N} \prod_{d\in\dirs} \mathbb{1}[\hat c_{i,d}=c_{i,d} \wedge \hat p_{i,d}=p_{i,d}], \\
  \TSS &= \tfrac{1}{3}\left(\acccolor + \accpass + \EM\right),
\end{align}
and, pooling over frames and directions,
\begin{equation}
  \FG = P(\hat p = \text{go} \mid p = \text{stop}), \qquad \FS = P(\hat p = \text{stop} \mid p = \text{go}).
\end{equation}

The Traffic Signal Score (TSS) is a convenient single number, but it has properties worth making
explicit:
\begin{itemize}
  \item \textbf{Trivial directions inflate it.} Because right is \emph{none} in $\sim$97\% of frames, a
        model that never predicts anything else for right gains almost a third of $\acccolor$ and
        $\accpass$ for free. For this reason, out-of-distribution results are also reported without the
        right direction.
  \item \textbf{It is symmetric in the errors that matter asymmetrically.} Predicting \emph{go} for a
        red signal (False-Go) can cause a collision; predicting \emph{stop} for a green one (False-Stop)
        causes an unnecessary stop and possibly a rear-end risk, but is generally the safer failure.
        TSS weights both equally. The False-Go and False-Stop rates are therefore reported alongside.
  \item \textbf{Pooling hides directions.} False-Go is pooled over directions, so the rare
        right-turn stops barely influence it even if the model misses all of them. Per-direction
        results are necessary to see this.
  \item \textbf{Accuracy and passability are coupled.} Under \cref{eq:pass-rule}, $\accpass\ge\acccolor$
        whenever the predicted passability follows the predicted color, because some color errors
        (green vs yellow vs none) do not change passability.
\end{itemize}
To make the per-direction numbers interpretable, each is compared with a \emph{majority} baseline: the
accuracy of always predicting that direction's most frequent color.

\section{Two design paradigms}
\label{sec:prob-paradigms}

Two ways of solving the task are compared in this thesis.

\paragraph{End-to-end (direction-aware).} The network maps the image directly to the 3$\times$2
outputs. Relevance, association and traffic rules are learned implicitly from data. An auxiliary
detector supplies localized evidence, but its output is consumed as features, not as decisions. The
advantages are flexibility (context such as lane markings or road layout can be used without being
modeled) and a single training objective. The drawbacks are opacity---it is hard to know \emph{why} a
maneuver was declared free---and the fact that the traffic rules are baked into the weights: a change
of convention requires relabeling and retraining.

\paragraph{Structured (detector + rules).} A detector predicts, for each light, a class that encodes
both its color and the set of ego maneuvers it controls (for instance \code{left+straight|red}, or
\code{irrelevant|green} for lights of other lanes or junctions). A fixed rule then assigns colors to
maneuvers with the same specificity hierarchy used to build the labels. When the detector is perfect,
this pipeline reproduces every label exactly (an ``oracle'' property verified in
\cref{sec:method-rfdetr}), so all its errors are detection, classification and relevance errors. The
advantages are interpretability, the possibility of changing the rule at inference time (US versus
German right turns) and the reuse of mature detection architectures. The drawbacks are the larger,
sparser class space and the dependence on a score threshold.

The empirical comparison of the two paradigms is the subject of \cref{sec:res-rfdetr}.
